\documentclass[11pt]{article}

% Practice problems for Lecture 1
\usepackage[margin=1in]{geometry}
\usepackage[T1]{fontenc}
\usepackage{lmodern}
\usepackage{amsmath,amssymb,amsthm}
\usepackage{enumitem}
\usepackage{hyperref}
\usepackage{xcolor}

\title{Lecture 1: Practice Problems}
\author{Linear Algebra}
\date{\today}

\theoremstyle{definition}
\newtheorem{problem}{Problem}[section]

\DeclareMathOperator{\im}{im}
\DeclareMathOperator{\rank}{rank}
\DeclareMathOperator{\Null}{null}
\DeclareMathOperator{\col}{col}
\DeclareMathOperator{\row}{row}
\DeclareMathOperator{\Span}{span}
\DeclareMathOperator{\RREF}{RREF}
\DeclareMathOperator{\tr}{tr}
\DeclareMathOperator{\diag}{diag}

\begin{document}

\maketitle

\section{Number Sets and Fields}

\begin{problem}\label{pr:numbers}
    For each of the numbers
    \[\frac{7}{3}, \qquad \sqrt{49}, \qquad \sqrt{12}, \qquad 0.25, \qquad \pi - \pi, \qquad (1 + i)(1 - i), \qquad 2 + 3i,\]
    decide which of the sets $\mathbb{Z}$, $\mathbb{Q}$, $\mathbb{R}$ and $\mathbb{C}$ it belongs to.
\end{problem}

\begin{problem}\label{pr:field-fail}
    Each of the following sets, with the usual addition and multiplication, is not a field. For each set, find a field axiom that fails.
    \begin{enumerate}[label=(\alph*)]
        \item $\mathbb{N} = \{0, 1, 2, \ldots\}$;
        \item the set $\mathbb{R}_{>0}$ of positive real numbers;
        \item the set $\{0, 1\}$;
        \item the set $2\mathbb{Z} = \{\ldots, -4, -2, 0, 2, 4, \ldots\}$ of even integers.
    \end{enumerate}
\end{problem}

\begin{problem}\label{pr:Qsqrt2}
    \begin{enumerate}[label=(\alph*)]
        \item Show that $\mathbb{Q}(\sqrt{2}) = \{a + b\sqrt{2} : a, b \in \mathbb{Q}\}$, with the usual addition and multiplication, is a field. (Hint: to find the inverse of $a + b\sqrt{2} \neq 0$, multiply by $\frac{a - b\sqrt{2}}{a - b\sqrt{2}}$. Why is $a^2 - 2b^2 \neq 0$?)
        \item Is $\{a + b\sqrt{2} : a, b \in \mathbb{Z}\}$ a field?
    \end{enumerate}
\end{problem}

\section{Vector Spaces}

\begin{problem}\label{pr:weird-ops}
    On $\mathbb{R}^2$, keep the usual vector addition, but define a new scalar multiplication. In each case, decide which vector space axioms hold and which fail.
    \begin{enumerate}[label=(\alph*)]
        \item $c \odot (x, y) = (cx, 0)$;
        \item $c \odot (x, y) = (cx, y)$.
    \end{enumerate}
\end{problem}

\begin{problem}\label{pr:positive-reals}
    On the set $V = \mathbb{R}_{>0}$ of positive real numbers, define ``addition'' and ``scalar multiplication'' by
    \[u \oplus v = uv \qquad \text{and} \qquad c \odot u = u^c \qquad (u, v \in V, \ c \in \mathbb{R}).\]
    Show that $V$ is a vector space over $\mathbb{R}$. What is its zero vector, and what is the additive inverse of $u$?
\end{problem}

\begin{problem}\label{pr:functions}
    \begin{enumerate}[label=(\alph*)]
        \item Show that the set of all functions $f: \mathbb{R} \to \mathbb{R}$, with $(f + g)(x) = f(x) + g(x)$ and $(cf)(x) = c f(x)$, is a vector space over $\mathbb{R}$. What is its zero vector?
        \item Is the set of all functions $f: \mathbb{R} \to \mathbb{R}$ with $f(0) = 1$, with the same operations, a vector space?
    \end{enumerate}
\end{problem}

\section{Bases and Dimension}

\begin{problem}\label{pr:indep}
    Decide whether the vectors $(1, 2, 1)$, $(2, 1, 0)$ and $(1, -1, -1)$ in $\mathbb{R}^3$ are linearly independent. If they are not, find a nontrivial linear combination that equals $0$, and find a basis and the dimension of their span.
\end{problem}

\begin{problem}\label{pr:basisU}
    One can check that $U = \{(x, y, z, w) \in \mathbb{R}^4 : x + y + z + w = 0 \text{ and } x = w\}$ is a vector space. Find a basis of $U$ and its dimension.
\end{problem}

\begin{problem}\label{pr:taylor-basis}
    Show that $\mathcal{B} = (1, \ x - 1, \ (x - 1)^2)$ is an ordered basis of the vector space $\mathbb{R}[x]_{\leq 2}$ of polynomials of degree at most $2$. Find the coordinate vectors $[x^2]_{\mathcal{B}}$ and $[3 + 2x + x^2]_{\mathcal{B}}$.
\end{problem}

\section{Linear Transformations}

\begin{problem}\label{pr:linear}
    Which of the following maps are linear transformations? Justify your answers.
    \begin{enumerate}[label=(\alph*)]
        \item $T: \mathbb{R}^2 \to \mathbb{R}^2$, $T(x, y) = (x - y, 3x)$;
        \item $T: \mathbb{R}^2 \to \mathbb{R}^2$, $T(x, y) = (x + y, 1)$;
        \item $T: \mathbb{R}^3 \to \mathbb{R}^2$, $T(x, y, z) = (xy, z)$;
        \item $D: \mathbb{R}[x]_{\leq 3} \to \mathbb{R}[x]_{\leq 3}$, $D(p) = p'$ (the derivative);
        \item $T: \mathbb{R} \to \mathbb{R}$, $T(x) = |x|$;
        \item $T: \mathbb{R}^2 \to \mathbb{R}^2$, $T(x, y) = (y, x)$.
    \end{enumerate}
\end{problem}

\begin{problem}\label{pr:kerim}
    Let $T: \mathbb{R}^3 \to \mathbb{R}^3$ be defined by $T(x, y, z) = (x + y, \ y + z, \ x - z)$. Find the kernel and the image of $T$. Is $T$ injective? Is it surjective?
\end{problem}

\begin{problem}\label{pr:gallery}
    For each of the following linear transformations of the plane, compute $T(e_1)$ and $T(e_2)$, sketch the image of the unit square and of the letter F from the gallery figure in the lecture, and find the kernel and the image of $T$.
    \begin{enumerate}[label=(\alph*)]
        \item the reflection in the line $y = x$, $T(x, y) = (y, x)$;
        \item the projection onto the $x$-axis, $T(x, y) = (x, 0)$;
        \item the shear $T(x, y) = (x + 2y, y)$;
        \item $T(x, y) = (-x, -y)$.
    \end{enumerate}
    In which cases is the image of the letter F mirrored?
\end{problem}

\begin{problem}\label{pr:derivative}
    Let $D: \mathbb{R}[x]_{\leq 3} \to \mathbb{R}[x]_{\leq 3}$, $D(p) = p'$. Find the kernel and the image of $D$, and decide whether $D$ is injective or surjective. What is $D \circ D \circ D \circ D$?
\end{problem}

\begin{problem}[$\star$]\label{pr:composition}
    Let $S$, $T$, $U$ be nonempty sets, and let $f: S \to T$ and $g: T \to U$ be functions.
    \begin{enumerate}[label=(\alph*)]
        \item Show that if $g \circ f$ is injective, then $f$ is injective.
        \item Give an example in which $g \circ f$ is injective, but $g$ is not.
        \item Show that if $g \circ f$ is surjective, then $g$ is surjective.
        \item Let $A: \mathbb{R}^m \to \mathbb{R}^n$ and $B: \mathbb{R}^n \to \mathbb{R}^m$ be linear transformations with $B \circ A = \mathrm{id}_{\mathbb{R}^m}$. Which of $A$ and $B$ must be injective, and which must be surjective?
    \end{enumerate}
\end{problem}

\begin{problem}[$\star$]\label{pr:distributive}
    Let $\mathcal{F}$ be the set of all functions $\mathbb{R} \to \mathbb{R}$, with pointwise addition $(f + g)(x) = f(x) + g(x)$ and composition $(f \circ g)(x) = f(g(x))$.
    \begin{enumerate}[label=(\alph*)]
        \item Show that $\circ$ is not commutative.
        \item Show that the right distributive law $(g + h) \circ f = (g \circ f) + (h \circ f)$ always holds.
        \item Show that the left distributive law $f \circ (g + h) = (f \circ g) + (f \circ h)$ can fail. (Try $f(x) = x^2$.)
        \item Show that the left distributive law holds whenever $f$ is linear, i.e., $f(x) = cx$ for some $c \in \mathbb{R}$. Compare with the distributive laws for matrix multiplication.
    \end{enumerate}
\end{problem}

\begin{problem}\label{pr:iso}
    \begin{enumerate}[label=(\alph*)]
        \item Show that $T: \mathbb{R}[x]_{\leq 3} \to \mathbb{R}^4$, $T(p) = \big(p(0), p'(0), \frac{1}{2}p''(0), \frac{1}{6}p'''(0)\big)$, is an isomorphism. (Hint: compute $T(a_0 + a_1 x + a_2 x^2 + a_3 x^3)$.)
        \item Show that the vector space $V$ of Problem~\ref{pr:positive-reals} is isomorphic to $\mathbb{R}$. (Hint: logarithm.)
        \item Which of the following pairs of real vector spaces are isomorphic: $\mathbb{R}[x]_{\leq 5}$ and $\mathbb{R}^6$; $\mathbb{C}$ (over $\mathbb{R}$) and $\mathbb{R}^3$; $\mathbb{R}^2 \times \mathbb{R}^3$ and $\mathbb{R}^5$?
    \end{enumerate}
\end{problem}

\section{Subspaces}

\begin{problem}\label{pr:subspaces}
    Which of the following subsets of $\mathbb{R}^3$ are subspaces?
    \begin{enumerate}[label=(\alph*)]
        \item $\{(x, y, z) : x = 2y\}$;
        \item $\{(x, y, z) : z = x^2\}$;
        \item $\{(x, y, z) : x + y = z + 1\}$;
        \item $\{(x, y, z) : x = 0 \text{ or } y = 0\}$;
        \item $\{(t, 2t, 3t) : t \in \mathbb{R}\}$;
        \item $\{(x, y, z) : x + y + z = 0 \text{ and } x = y\}$.
    \end{enumerate}
\end{problem}

\begin{problem}\label{pr:intersection}
    Let $U$ and $W$ be subspaces of a vector space $V$.
    \begin{enumerate}[label=(\alph*)]
        \item Show that $U \cap W$ is a subspace of $V$.
        \item Show that $U \cup W$ need not be a subspace.
        \item Show that $U \cup W$ is a subspace if and only if $U \subseteq W$ or $W \subseteq U$.
    \end{enumerate}
\end{problem}

\begin{problem}\label{pr:poly-subspace}
    Show that $\{p \in \mathbb{R}[x]_{\leq 3} : p(1) = 0\}$ is a subspace of $\mathbb{R}[x]_{\leq 3}$ by writing it as the kernel of a linear transformation. Find a basis of this subspace and its dimension.
\end{problem}

\section{Matrices}

\subsection{Matrix Operations}

\begin{problem}\label{pr:matvec}
    Let $A = \begin{pmatrix} 1 & 0 & 2 \\ -1 & 3 & 1 \end{pmatrix}$ and $x = (2, 1, -1)^\top$.
    \begin{enumerate}[label=(\alph*)]
        \item Compute $Ax$ in two ways: entry by entry, and as a linear combination of the columns of $A$.
        \item Compute $A^\top$, $AA^\top$ and $A^\top A$. Which of these matrices are symmetric?
    \end{enumerate}
\end{problem}

\begin{problem}\label{pr:sizes}
    Let $A$ be a $2 \times 3$ matrix, $B$ a $3 \times 3$ matrix and $C$ a $3 \times 2$ matrix. Which of the following expressions are defined, and what are their sizes?
    \[AB, \qquad BA, \qquad AC, \qquad CA, \qquad A + C^\top, \qquad A^\top + C, \qquad ABC.\]
\end{problem}

\begin{problem}\label{pr:reflect-rotate}
    Let $R = \begin{pmatrix} 0 & -1 \\ 1 & 0 \end{pmatrix}$ be the rotation by $90^\circ$, and $F = \begin{pmatrix} 1 & 0 \\ 0 & -1 \end{pmatrix}$ the reflection in the $x$-axis. Compute $FR$ and $RF$, and describe both as geometric transformations. What does this tell you about the order in which linear transformations are composed?
\end{problem}

\begin{problem}\label{pr:rotation-addition}
    Let $R_\theta = \begin{pmatrix} \cos\theta & -\sin\theta \\ \sin\theta & \cos\theta \end{pmatrix}$ be the rotation by the angle $\theta$.
    \begin{enumerate}[label=(\alph*)]
        \item Explain geometrically why $R_\alpha R_\beta = R_{\alpha + \beta}$.
        \item Multiply out $R_\alpha R_\beta$, and deduce the addition formulas for $\cos(\alpha + \beta)$ and $\sin(\alpha + \beta)$.
        \item Show that $(R_\theta)^n = R_{n\theta}$, and conclude that $(R_{2\pi/n})^n = I_2$.
    \end{enumerate}
\end{problem}

\begin{problem}[$\star$]\label{pr:complex-matrix}
    For a complex number $z = a + bi$, let $M(z) = \begin{pmatrix} a & -b \\ b & a \end{pmatrix}$.
    \begin{enumerate}[label=(\alph*)]
        \item Show that $M(z + w) = M(z) + M(w)$ and $M(zw) = M(z) M(w)$ for all $z, w \in \mathbb{C}$.
        \item Show that $M(1 + i) = \sqrt{2}\, R_{\pi/4}$, with $R_\theta$ as in Problem~\ref{pr:rotation-addition}.
        \item Use parts (a) and (b) to compute $(1 + i)^{10}$.
    \end{enumerate}
\end{problem}

\subsection{Row Reduction}

\begin{problem}\label{pr:ref}
    Which of the following matrices are in row echelon form? Which are in reduced row echelon form?
    \[\begin{pmatrix} 1 & 3 & 0 & 2 \\ 0 & 0 & 1 & 4 \\ 0 & 0 & 0 & 0 \end{pmatrix}, \quad \begin{pmatrix} 1 & 3 & 1 \\ 0 & 1 & 2 \\ 0 & 0 & 0 \end{pmatrix}, \quad \begin{pmatrix} 0 & 0 & 1 \\ 0 & 1 & 0 \\ 1 & 0 & 0 \end{pmatrix}, \quad \begin{pmatrix} 1 & 0 & 0 \\ 0 & 0 & 0 \\ 0 & 0 & 1 \end{pmatrix}, \quad \begin{pmatrix} 2 & 1 \\ 0 & 1 \end{pmatrix}.\]
\end{problem}

\begin{problem}\label{pr:rref-compute}
    Find the reduced row echelon form and the rank of
    \[P = \begin{pmatrix} 1 & 2 & 1 & 0 \\ 2 & 4 & 0 & 2 \\ 1 & 2 & 3 & -2 \end{pmatrix}.\]
\end{problem}

\begin{problem}[$\star$]\label{pr:system}
    Consider the system of linear equations
    \begin{align*}
        w + 3x + 2y + 2z &= 0, \\
        w + 4x + y &= 0, \\
        3w + 5x + 10y + 14z &= 0, \\
        2w + 5x + 5y + 6z &= 0.
    \end{align*}
    \begin{enumerate}[label=(\alph*)]
        \item Find all solutions $(w, x, y, z) \in \mathbb{R}^4$ by bringing the coefficient matrix into reduced row echelon form.
        \item Check that $(-5, 1, 1, 0)$ and $(0, 2, -8, 5)$ are solutions. Is every solution a scalar multiple of $(-5, 1, 1, 0)$?
        \item What is the dimension of the space of solutions?
    \end{enumerate}
\end{problem}

\subsection{The Inverse of a Matrix}

\begin{problem}\label{pr:inverse2}
    Find the inverse of $A = \begin{pmatrix} 3 & 1 \\ 5 & 2 \end{pmatrix}$ by solving the equation $Ax = y$ for $x$, as in the lecture. Check your answer by computing $AA^{-1}$.
\end{problem}

\begin{problem}\label{pr:inverse3}
    Use the Gauss--Jordan method to find the inverse of
    \[A = \begin{pmatrix} 1 & 2 & 1 \\ 0 & 1 & 1 \\ 1 & 1 & 1 \end{pmatrix},\]
    and use it to solve the system $Ax = (1, 2, 3)^\top$.
\end{problem}

\begin{problem}\label{pr:arithmetic}
    \begin{enumerate}[label=(\alph*)]
        \item Apply the Gauss--Jordan method to $A = \begin{pmatrix} 1 & 2 & 3 \\ 2 & 3 & 4 \\ 3 & 4 & 5 \end{pmatrix}$ to show that $A$ is singular, and find a nonzero vector $x$ with $Ax = 0$.
        \item Show more generally that every $3 \times 3$ matrix whose rows form an arithmetic progression (i.e., row $2$ $-$ row $1$ $=$ row $3$ $-$ row $2$) is singular.
    \end{enumerate}
\end{problem}

\begin{problem}\label{pr:elementary}
    Let $A = \begin{pmatrix} 0 & 2 \\ 1 & 3 \end{pmatrix}$. Find elementary matrices $E_1$, $E_2$, $E_3$ such that $E_3 E_2 E_1 A = I_2$. Use them to write $A^{-1}$, and also $A$ itself, as a product of elementary matrices.
\end{problem}

\begin{problem}[$\star$]\label{pr:idempotent}
    A square matrix $A$ with $A^2 = A$ is called idempotent.
    \begin{enumerate}[label=(\alph*)]
        \item Show that the only invertible idempotent matrix is the identity.
        \item Show that $P = \begin{pmatrix} 1 & 1 \\ 0 & 0 \end{pmatrix}$ is idempotent, and find the kernel and the image of $T_P$. Describe $T_P$ geometrically.
        \item Show that if $A$ is idempotent, then so is $I - A$, and that $A(I - A) = 0$.
        \item Let $A$ be an idempotent $n \times n$ matrix. Show that $\im(T_A) \cap \ker(T_A) = \{0\}$, and that every $v \in \mathbb{R}^n$ can be written as $v = Av + (v - Av)$, where $Av \in \im(T_A)$ and $v - Av \in \ker(T_A)$.
    \end{enumerate}
\end{problem}

\begin{problem}\label{pr:nilpotent}
    Let $N = \begin{pmatrix} 0 & 1 & 2 \\ 0 & 0 & 3 \\ 0 & 0 & 0 \end{pmatrix}$.
    \begin{enumerate}[label=(\alph*)]
        \item Compute $N^2$ and $N^3$.
        \item Show that $(I - N)(I + N + N^2) = I$, and compute $(I - N)^{-1}$.
        \item More generally, show that if $N^k = 0$ for some $k \geq 1$, then $I - N$ is invertible, with $(I - N)^{-1} = I + N + N^2 + \cdots + N^{k-1}$. (Compare with the geometric series $\frac{1}{1 - x} = 1 + x + x^2 + \cdots$.)
    \end{enumerate}
\end{problem}

\begin{problem}\label{pr:interpolation}
    Let $T: \mathbb{R}[x]_{\leq 2} \to \mathbb{R}^3$, $T(p) = \big(p(-1), p(0), p(1)\big)$.
    \begin{enumerate}[label=(\alph*)]
        \item Find the matrix of $T$ with respect to the ordered basis $(1, x, x^2)$ of $\mathbb{R}[x]_{\leq 2}$ and the standard basis of $\mathbb{R}^3$.
        \item Show that this matrix is invertible, and conclude that $T$ is an isomorphism.
        \item Find the polynomial $p \in \mathbb{R}[x]_{\leq 2}$ with $p(-1) = 0$, $p(0) = 1$, $p(1) = 0$, and the polynomial $q$ with $q(-1) = 1$, $q(0) = 2$, $q(1) = 5$.
    \end{enumerate}
\end{problem}

\subsection{Special Types of Matrices}

\begin{problem}\label{pr:dims-special}
    Find the dimensions of the following vector spaces, and a basis of each of the first two:
    \begin{enumerate}[label=(\alph*)]
        \item $\mathrm{Sym}_2$;
        \item $\mathrm{Skew}_2$;
        \item $\mathbb{R}^{2 \times 3}$;
        \item $\mathrm{Sym}_4$;
        \item $\mathrm{Skew}_4$;
        \item the space of upper triangular $3 \times 3$ matrices with trace $0$.
    \end{enumerate}
\end{problem}

\begin{problem}\label{pr:groups}
    Which of the following sets of matrices are matrix groups, i.e., contain the identity and are closed under multiplication and under taking inverses?
    \begin{enumerate}[label=(\alph*)]
        \item the invertible diagonal $n \times n$ matrices;
        \item the invertible symmetric $2 \times 2$ matrices;
        \item the $2 \times 2$ matrices $A$ with $A^2 = I_2$;
        \item the upper triangular $n \times n$ matrices whose diagonal entries are all equal to $1$;
        \item the $n \times n$ permutation matrices. How many of them are there?
    \end{enumerate}
\end{problem}

\begin{problem}\label{pr:inverse-triangular}
    Use the Gauss--Jordan method to find the inverse of $\begin{pmatrix} 1 & 2 & 3 \\ 0 & 1 & 4 \\ 0 & 0 & 1 \end{pmatrix}$. Then explain why the inverse of every invertible upper triangular matrix is again upper triangular.
\end{problem}

\begin{problem}\label{pr:orthogonal2}
    Find all orthogonal $2 \times 2$ matrices. (Hint: for $Q = \begin{pmatrix} a & b \\ c & d \end{pmatrix}$, the condition $Q^\top Q = I_2$ means $a^2 + c^2 = 1$, $b^2 + d^2 = 1$ and $ab + cd = 0$. Write $(a, c) = (\cos\theta, \sin\theta)$.) Which of them are rotations, and which are reflections?
\end{problem}

\subsection{Trace and Determinant}

\begin{problem}[$\star$]\label{pr:tridiagonal}
    Let $D_n$ be the determinant of the $n \times n$ matrix with $2$s on the diagonal, $1$s directly above and below the diagonal, and $0$s everywhere else; for example,
    \[D_3 = \det \begin{pmatrix} 2 & 1 & 0 \\ 1 & 2 & 1 \\ 0 & 1 & 2 \end{pmatrix}.\]
    \begin{enumerate}[label=(\alph*)]
        \item Compute $D_1$, $D_2$ and $D_3$.
        \item Expand along the first row to show that $D_n = 2D_{n-1} - D_{n-2}$ for $n \geq 3$.
        \item Find a formula for $D_n$.
    \end{enumerate}
\end{problem}

\begin{problem}\label{pr:det4}
    Compute the determinant of
    \[A = \begin{pmatrix} 2 & 0 & 1 & 3 \\ 1 & 1 & 0 & 2 \\ 0 & 3 & 1 & 1 \\ 1 & 0 & 2 & 1 \end{pmatrix}\]
    in two ways: by cofactor expansion along a well-chosen row or column, and by row reduction. Is $A$ invertible?
\end{problem}

\begin{problem}\label{pr:det-rules}
    Let $A$ and $B$ be $3 \times 3$ matrices with $\det(A) = 2$ and $\det(B) = -3$. Compute
    \[\det(AB), \qquad \det(A^\top B^2), \qquad \det(3A), \qquad \det(A^{-1}B), \qquad \det(BAB^{-1}).\]
    Let $C$ be obtained from $A$ by swapping rows $1$ and $3$ and then adding $5$ times row $1$ to row $2$. What is $\det(C)$?
\end{problem}

\begin{problem}\label{pr:vandermonde}
    \begin{enumerate}[label=(\alph*)]
        \item Use row operations to show that
        \[\det \begin{pmatrix} 1 & a & a^2 \\ 1 & b & b^2 \\ 1 & c & c^2 \end{pmatrix} = (b - a)(c - a)(c - b).\]
        \item Conclude that for distinct numbers $a$, $b$, $c$ and arbitrary values $y_1$, $y_2$, $y_3$, there is exactly one polynomial $p \in \mathbb{R}[x]_{\leq 2}$ with $p(a) = y_1$, $p(b) = y_2$ and $p(c) = y_3$. Compare with Problem~\ref{pr:interpolation}.
    \end{enumerate}
\end{problem}

\begin{problem}[$\star$]\label{pr:lambda}
    Let $A = \begin{pmatrix} 3 & 5 & 3 \\ 1 & 7 & 3 \\ 1 & 2 & 8 \end{pmatrix}$.
    \begin{enumerate}[label=(\alph*)]
        \item For which of the numbers $\lambda = 2, 3, 5$ is the matrix $A - \lambda I_3$ singular?
        \item For each such $\lambda$, find a nonzero vector $x$ with $Ax = \lambda x$.
        \item Without computing any determinant, find another number $\lambda$ and a nonzero vector $x$ with $Ax = \lambda x$. (Hint: add up the entries in each row of $A$.)
    \end{enumerate}
\end{problem}

\begin{problem}[$\star$]\label{pr:upper5}
    Let
    \[A = \begin{pmatrix} 1 & 2 & 3 & 4 & 5 \\ 0 & 2 & 3 & 4 & 5 \\ 0 & 0 & 3 & 4 & 5 \\ 0 & 0 & 0 & 4 & 5 \\ 0 & 0 & 0 & 0 & 5 \end{pmatrix}.\]
    \begin{enumerate}[label=(\alph*)]
        \item Compute $\det(A)$. Is $A$ invertible? Can $A$ be transformed into $I_5$ by elementary row operations?
        \item Find the last row of $A^2$ without computing all of $A^2$.
        \item Explain why $A - I_5$ and $A - 2I_5$ are singular. Find nonzero vectors $x$ and $y$ with $Ax = x$ and $Ay = 2y$.
    \end{enumerate}
\end{problem}

\begin{problem}\label{pr:area}
    \begin{enumerate}[label=(\alph*)]
        \item Find the area of the triangle with vertices $(1, 1)$, $(4, 2)$ and $(2, 5)$. (Hint: move one vertex to the origin.)
        \item Let $T(x, y) = (2x - y, x + 3y)$. Find the area of the image under $T$ of the unit square and of the unit disk.
        \item Show that the image of the unit disk under $T(x, y) = (ax, by)$, with $a, b > 0$, is the ellipse $\frac{x^2}{a^2} + \frac{y^2}{b^2} \leq 1$, and conclude that its area is $\pi ab$.
        \item Which of the transformations of Problem~\ref{pr:gallery} reverse the orientation? Compare with your sketches of the letter F.
    \end{enumerate}
\end{problem}

\begin{problem}\label{pr:fibonacci}
    The Fibonacci numbers are defined by $F_0 = 0$, $F_1 = 1$ and $F_{n+1} = F_n + F_{n-1}$. Let $M = \begin{pmatrix} 1 & 1 \\ 1 & 0 \end{pmatrix}$.
    \begin{enumerate}[label=(\alph*)]
        \item Show by induction that $M^n = \begin{pmatrix} F_{n+1} & F_n \\ F_n & F_{n-1} \end{pmatrix}$ for $n \geq 1$, and compute $M^{10}$.
        \item Use the product rule for determinants to prove Cassini's identity $F_{n+1}F_{n-1} - F_n^2 = (-1)^n$. Check it for $n = 10$.
    \end{enumerate}
\end{problem}

\begin{problem}\label{pr:trace}
    \begin{enumerate}[label=(\alph*)]
        \item Find $2 \times 2$ matrices $A$ and $B$ with $AB - BA = \begin{pmatrix} 1 & 0 \\ 0 & -1 \end{pmatrix}$. Why is there no such pair with $AB - BA = I_2$?
        \item Show that if $\tr(A^\top A) = 0$, then $A = 0$.
        \item Let $A$ be an $n \times n$ matrix with $\tr(AX) = 0$ for every $n \times n$ matrix $X$. Show that $A = 0$. (Hint: try $X = e_{ji}$, the matrix with a single $1$ in position $(j, i)$.)
    \end{enumerate}
\end{problem}

\begin{problem}[$\star$]\label{pr:01-matrices}
    \begin{enumerate}[label=(\alph*)]
        \item How many of the $16$ matrices $\begin{pmatrix} a & b \\ c & d \end{pmatrix}$ with all entries in $\{0, 1\}$ are invertible?
        \item Show that a square matrix $A$ with integer entries has an inverse with integer entries if and only if $\det(A) = \pm 1$.
    \end{enumerate}
\end{problem}

\section{Null Space, Column Space and Rank}

\begin{problem}\label{pr:bases-B}
    Let $B = \begin{pmatrix} 1 & 2 & 0 & 3 \\ 2 & 4 & 1 & 8 \\ 1 & 2 & 1 & 5 \end{pmatrix}$. Find $\RREF(B)$, the rank of $B$, and bases of $\Null(B)$, $\col(B)$ and $\row(B)$. Check that your answer agrees with the rank--nullity theorem.
\end{problem}

\begin{problem}[$\star$]\label{pr:onto}
    \begin{enumerate}[label=(\alph*)]
        \item Let $T: \mathbb{R}^5 \to \mathbb{R}^3$ be a surjective linear transformation. What is the dimension of $\ker(T)$?
        \item Can a linear transformation $\mathbb{R}^3 \to \mathbb{R}^5$ be surjective? Can it be injective?
        \item Can a $4 \times 6$ matrix $A$ have $\Null(A) = \{0\}$?
        \item Let $A$ be a $5 \times 3$ matrix such that $Ax = 0$ only for $x = 0$. What is the rank of $A$?
    \end{enumerate}
\end{problem}

\begin{problem}[$\star$]\label{pr:intersect}
    \begin{enumerate}[label=(\alph*)]
        \item Let $V$ and $W$ be $4$-dimensional subspaces of $\mathbb{R}^7$. Show that $V \cap W \neq \{0\}$. (Hint: put a basis of $V$ and a basis of $W$ as the columns of a $7 \times 8$ matrix, and apply the rank--nullity theorem.)
        \item More generally, show that if $V$ and $W$ are subspaces of $\mathbb{R}^n$ with $\dim V + \dim W > n$, then $V \cap W \neq \{0\}$.
        \item Find two $3$-dimensional subspaces of $\mathbb{R}^6$ whose intersection is $\{0\}$.
    \end{enumerate}
\end{problem}

\begin{problem}\label{pr:rankAB}
    \begin{enumerate}[label=(\alph*)]
        \item Show that $\rank(AB) \leq \rank(A)$ and $\rank(AB) \leq \rank(B)$ whenever $AB$ is defined. (Hint: $\col(AB) \subseteq \col(A)$ and $\Null(B) \subseteq \Null(AB)$.)
        \item Conclude that if $A$ is $3 \times 2$ and $B$ is $2 \times 3$, then the $3 \times 3$ matrix $AB$ is never invertible.
    \end{enumerate}
\end{problem}

\begin{problem}\label{pr:ATA}
    Let $A$ be an $m \times n$ matrix.
    \begin{enumerate}[label=(\alph*)]
        \item Show that $\Null(A^\top A) = \Null(A)$. (Hint: if $A^\top A x = 0$, then $x^\top A^\top A x = 0$; what is $x^\top A^\top A x$ in terms of the entries of $Ax$?)
        \item Conclude that $\rank(A^\top A) = \rank(A)$.
        \item Show that if the columns of $A$ are linearly independent, then $A^\top A$ is invertible. (This is why the normal equations $(A^\top A)w = A^\top y$ of linear regression have a unique solution.)
    \end{enumerate}
\end{problem}

\section{Eigenvalues and Eigenvectors}

\begin{problem}\label{pr:eig2}
    Let $A = \begin{pmatrix} 1 & 2 \\ 3 & 2 \end{pmatrix}$.
    \begin{enumerate}[label=(\alph*)]
        \item Compute the characteristic polynomial of $A$, its eigenvalues, and its eigenspaces.
        \item Check that the eigenvalues add up to $\tr(A)$ and multiply to $\det(A)$.
        \item Diagonalize $A$, i.e., find $P$ and $D$ with $A = PDP^{-1}$.
        \item Find a formula for $A^k$, and check it for $k = 1$.
    \end{enumerate}
\end{problem}

\begin{problem}\label{pr:eig-triangular}
    Find the eigenvalues and eigenspaces of $L = \begin{pmatrix} 2 & 0 & 0 \\ 1 & 3 & 0 \\ 4 & -1 & 1 \end{pmatrix}$. Is $L$ diagonalizable?
\end{problem}

\begin{problem}\label{pr:eig-defective}
    Let $N = \begin{pmatrix} 2 & 1 & 0 \\ 0 & 2 & 0 \\ 0 & 0 & 3 \end{pmatrix}$. Find the algebraic and geometric multiplicity of each eigenvalue of $N$, and decide whether $N$ is diagonalizable.
\end{problem}

\begin{problem}\label{pr:eig-geometric}
    Without computing any characteristic polynomial, use geometric intuition to find the eigenvalues and eigenspaces of the following linear transformations of the plane. Then check your answers with a computation.
    \begin{enumerate}[label=(\alph*)]
        \item the reflection in the $x$-axis;
        \item the orthogonal projection onto the line $y = x$, with matrix $\frac{1}{2}\begin{pmatrix} 1 & 1 \\ 1 & 1 \end{pmatrix}$;
        \item the shear $T(x, y) = (x + 2y, y)$;
        \item the rotation by $90^\circ$;
        \item $T(x, y) = (2x, 3y)$;
        \item $T(x, y) = (-x, -y)$.
    \end{enumerate}
    Which of these transformations are diagonalizable?
\end{problem}

\begin{problem}[$\star$]\label{pr:eig-gre}
    A real $3 \times 3$ matrix $A$ has the eigenvalues $-1$, $2$ and $4$.
    \begin{enumerate}[label=(\alph*)]
        \item Is $A$ diagonalizable? Is $A$ invertible?
        \item Compute $\tr(A)$, $\det(A)$, $\tr(A^2)$ and $\det(A^3)$.
        \item Find the eigenvalues of $A^{-1}$, $A^2$ and $A + 2I_3$.
    \end{enumerate}
\end{problem}

\begin{problem}[$\star$]\label{pr:eig-idempotent}
    Let $P$ be an $n \times n$ matrix with $P^2 = P$ (compare Problem~\ref{pr:idempotent}).
    \begin{enumerate}[label=(\alph*)]
        \item Show that every eigenvalue of $P$ is $0$ or $1$.
        \item Show that every nonzero vector of $\im(T_P)$ is an eigenvector for the eigenvalue $1$, and that every nonzero vector of $\ker(T_P)$ is an eigenvector for the eigenvalue $0$. Use Problem~\ref{pr:idempotent}(d) to conclude that $P$ is diagonalizable.
        \item Which of the following statements must be true: $P$ is invertible; $P$ is the identity or the zero matrix; $P$ is diagonalizable?
    \end{enumerate}
\end{problem}

\begin{problem}\label{pr:eig-props}
    Let $A$ and $B$ be $n \times n$ matrices.
    \begin{enumerate}[label=(\alph*)]
        \item Show that $A$ and $A^\top$ have the same eigenvalues, and give an example where they have different eigenvectors.
        \item Show that if $A$ is invertible, then $AB$ and $BA$ are similar, and hence have the same eigenvalues. (Hint: $BA = A^{-1}(AB)A$.)
        \item Show that every real eigenvalue of an orthogonal matrix $Q$ is $1$ or $-1$. Hint: if $Qv = \lambda v$, compute $(Qv)^\top (Qv)$ in two ways.
        \item Show that if $A^k = 0$ for some $k \geq 1$, then $0$ is the only eigenvalue of $A$.
        \item Show that $\begin{pmatrix} 1 & 1 \\ 0 & 1 \end{pmatrix}$ and $I_2$ have the same characteristic polynomial, but are not similar.
    \end{enumerate}
\end{problem}

\begin{problem}[$\star$]\label{pr:eig-recurrence}
    The sequence $a_0 = 0$, $a_1 = 1$, $a_{n+1} = a_n + 2a_{n-1}$ starts with $0, 1, 1, 3, 5, 11, \ldots$.
    \begin{enumerate}[label=(\alph*)]
        \item Find a $2 \times 2$ matrix $M$ with $(a_{n+1}, a_n)^\top = M (a_n, a_{n-1})^\top$.
        \item Find the eigenvalues and eigenvectors of $M$, and derive a formula for $a_n$, as for Binet's formula in the lecture.
        \item Compute $a_{10}$ with your formula.
    \end{enumerate}
\end{problem}

\begin{problem}\label{pr:eig-markov}
    Every day, a visitor of a website either only browses or makes a purchase. A browsing visitor browses again the next day with probability $0.7$, and a buying visitor buys again with probability $0.6$. With $x_k = (b_k, p_k)^\top$ the probabilities of browsing and buying on day $k$, we have $x_{k+1} = Ax_k$ with
    \[A = \begin{pmatrix} 0.7 & 0.4 \\ 0.3 & 0.6 \end{pmatrix}.\]
    \begin{enumerate}[label=(\alph*)]
        \item Find the eigenvalues and eigenvectors of $A$, and the stationary distribution.
        \item For $x_0 = (1, 0)^\top$, find a formula for $x_k$, and check it for $k = 2$. What happens as $k \to \infty$?
    \end{enumerate}
\end{problem}

\begin{problem}\label{pr:eig-power}
    Let $A = \begin{pmatrix} 2 & 1 \\ 1 & 2 \end{pmatrix}$ and $x_0 = (1, 0)^\top$.
    \begin{enumerate}[label=(\alph*)]
        \item Compute $Ax_0$, $A^2x_0$, $A^3x_0$ and $A^4x_0$. What happens to the ratio of the two entries?
        \item Explain your observation by writing $x_0$ as a linear combination of eigenvectors of $A$, and find a formula for $A^k x_0$.
    \end{enumerate}
\end{problem}

\begin{problem}[$\star$]\label{pr:eig-complex}
    \begin{enumerate}[label=(\alph*)]
        \item Show that the matrix $M(a + bi) = \begin{pmatrix} a & -b \\ b & a \end{pmatrix}$ of Problem~\ref{pr:complex-matrix} has the complex eigenvalues $a \pm bi$.
        \item Find the complex eigenvalues of $\begin{pmatrix} 3 & -4 \\ 4 & 3 \end{pmatrix}$, and write this matrix as $r R_\theta$ with $r > 0$ and a rotation $R_\theta$. How are $r$ and $\theta$ related to the eigenvalues?
    \end{enumerate}
\end{problem}

\section*{Answers and Hints}
\addcontentsline{toc}{section}{Answers and Hints}

\begin{description}[leftmargin=0pt, itemsep=0.3em]
    \item[\ref{pr:numbers}.] $\frac{7}{3}, 0.25 \in \mathbb{Q}, \mathbb{R}, \mathbb{C}$; $\sqrt{49} = 7$, $\pi - \pi = 0$ and $(1 + i)(1 - i) = 2$ lie in all four sets; $\sqrt{12} = 2\sqrt{3} \in \mathbb{R}, \mathbb{C}$; $2 + 3i \in \mathbb{C}$ only.
    \item[\ref{pr:field-fail}.] (a) Additive inverses: $-1 \notin \mathbb{N}$. (b) There is no additive identity, since $0 \notin \mathbb{R}_{>0}$. (c) Not closed under addition: $1 + 1 = 2$. (d) There is no multiplicative identity, since $1 \notin 2\mathbb{Z}$.
    \item[\ref{pr:Qsqrt2}.] (a) $(a + b\sqrt{2})^{-1} = \frac{a}{a^2 - 2b^2} - \frac{b}{a^2 - 2b^2}\sqrt{2}$; if $a^2 = 2b^2$ with $(a, b) \neq (0, 0)$, then $\sqrt{2} = |a/b|$ would be rational. (b) No: $2$ has no inverse, since $\frac{1}{2} = a + b\sqrt{2}$ forces $b = 0$ and $a = \frac{1}{2} \notin \mathbb{Z}$.
    \item[\ref{pr:weird-ops}.] (a) Only $1 \odot v = v$ fails. (b) Only $(a + b) \odot v = a \odot v + b \odot v$ fails: the second coordinates are $y$ and $2y$.
    \item[\ref{pr:positive-reals}.] The zero vector is $1$, and the additive inverse of $u$ is $\frac{1}{u}$. For example, $(a + b) \odot u = u^{a + b} = u^a u^b = (a \odot u) \oplus (b \odot u)$.
    \item[\ref{pr:functions}.] (a) The zero vector is the zero function $x \mapsto 0$. (b) No: the zero function is not in the set (and the set is not closed under addition).
    \item[\ref{pr:indep}.] Dependent: $(1, 2, 1) - (2, 1, 0) + (1, -1, -1) = 0$. The span has basis $\{(1, 2, 1), (2, 1, 0)\}$ and dimension $2$.
    \item[\ref{pr:basisU}.] $\dim U = 2$, with basis $(1, 0, -2, 1)$, $(0, 1, -1, 0)$.
    \item[\ref{pr:taylor-basis}.] Substitute $x = t + 1$: $[x^2]_{\mathcal{B}} = (1, 2, 1)$ and $[3 + 2x + x^2]_{\mathcal{B}} = (6, 4, 1)$.
    \item[\ref{pr:linear}.] (a) Linear. (b) Not linear: $T(0, 0) \neq (0, 0)$. (c) Not linear: $T(2v) = 4T(v)$ in the first coordinate. (d) Linear. (e) Not linear: $T(-1) = 1 \neq -T(1)$. (f) Linear.
    \item[\ref{pr:kerim}.] $\ker(T) = \Span\{(1, -1, 1)\}$ and $\im(T) = \{(a, b, c) : c = a - b\} = \Span\{(1, 0, 1), (1, 1, 0)\}$. $T$ is neither injective nor surjective.
    \item[\ref{pr:gallery}.] (a) $T(e_1) = e_2$, $T(e_2) = e_1$; bijective; the F is mirrored. (b) $\ker(T)$ is the $y$-axis and $\im(T)$ is the $x$-axis; the square collapses onto the segment from $(0, 0)$ to $(1, 0)$. (c) Bijective; the square becomes the parallelogram with vertices $(0, 0), (1, 0), (3, 1), (2, 1)$. (d) Rotation by $180^\circ$; bijective; the F is not mirrored. Only (a) mirrors the F.
    \item[\ref{pr:derivative}.] $\ker(D)$ is the set of constant polynomials, and $\im(D) = \mathbb{R}[x]_{\leq 2}$. $D$ is neither injective nor surjective. $D \circ D \circ D \circ D = 0$.
    \item[\ref{pr:composition}.] (a) If $f(a) = f(b)$, then $g(f(a)) = g(f(b))$, so $a = b$. (b) $f: \{1\} \to \{1, 2\}$, $f(1) = 1$, and $g: \{1, 2\} \to \{1\}$. (d) $A$ is injective and $B$ is surjective.
    \item[\ref{pr:distributive}.] (a) For example, $f(x) = x + 1$ and $g(x) = 2x$. (c) With $f(x) = x^2$ and $g = h = \mathrm{id}$: $f(2x) = 4x^2 \neq 2x^2$. (d) $c(g(x) + h(x)) = cg(x) + ch(x)$.
    \item[\ref{pr:iso}.] (a) $T(a_0 + a_1 x + a_2 x^2 + a_3 x^3) = (a_0, a_1, a_2, a_3)$. (b) $\log: V \to \mathbb{R}$, since $\log(uv) = \log u + \log v$ and $\log(u^c) = c\log u$. (c) Yes; no (dimensions $2$ and $3$); yes.
    \item[\ref{pr:subspaces}.] (a) yes; (b) no; (c) no; (d) no; (e) yes; (f) yes.
    \item[\ref{pr:intersection}.] (b) The $x$-axis and the $y$-axis in $\mathbb{R}^2$: $(1, 0) + (0, 1) = (1, 1)$ lies in neither. (c) If $u \in U \setminus W$ and $w \in W \setminus U$, then $u + w$ lies in neither $U$ nor $W$.
    \item[\ref{pr:poly-subspace}.] It is the kernel of $p \mapsto p(1)$. A basis is $x - 1, \ x^2 - 1, \ x^3 - 1$, so the dimension is $3$.
    \item[\ref{pr:matvec}.] (a) $Ax = 0$. (b) $AA^\top = \begin{pmatrix} 5 & 1 \\ 1 & 11 \end{pmatrix}$ and $A^\top A = \begin{pmatrix} 2 & -3 & 1 \\ -3 & 9 & 3 \\ 1 & 3 & 5 \end{pmatrix}$, both symmetric.
    \item[\ref{pr:sizes}.] $AB$: $2 \times 3$; $BA$: undefined; $AC$: $2 \times 2$; $CA$: $3 \times 3$; $A + C^\top$: $2 \times 3$; $A^\top + C$: $3 \times 2$; $ABC$: $2 \times 2$.
    \item[\ref{pr:reflect-rotate}.] $FR = \begin{pmatrix} 0 & -1 \\ -1 & 0 \end{pmatrix}$ is the reflection in the line $y = -x$, and $RF = \begin{pmatrix} 0 & 1 \\ 1 & 0 \end{pmatrix}$ is the reflection in the line $y = x$. So $FR \neq RF$.
    \item[\ref{pr:rotation-addition}.] (b) $\cos(\alpha + \beta) = \cos\alpha\cos\beta - \sin\alpha\sin\beta$ and $\sin(\alpha + \beta) = \sin\alpha\cos\beta + \cos\alpha\sin\beta$.
    \item[\ref{pr:complex-matrix}.] (c) $M(1 + i)^{10} = 2^5 R_{10\pi/4} = 32 R_{\pi/2} = \begin{pmatrix} 0 & -32 \\ 32 & 0 \end{pmatrix} = M(32i)$, so $(1 + i)^{10} = 32i$.
    \item[\ref{pr:ref}.] The first is in RREF. The second is in REF, but not RREF. The third and fourth are not in REF (the pivots move left; a zero row lies above a nonzero row). The fifth is not in REF with our convention, since its first pivot is $2$.
    \item[\ref{pr:rref-compute}.] $\RREF(P) = \begin{pmatrix} 1 & 2 & 0 & 1 \\ 0 & 0 & 1 & -1 \\ 0 & 0 & 0 & 0 \end{pmatrix}$, and $\rank(P) = 2$.
    \item[\ref{pr:system}.] The RREF has rows $(1, 0, 5, 8)$ and $(0, 1, -1, -2)$, so the solutions are
    \[(w, x, y, z) = y\,(-5, 1, 1, 0) + z\,(-8, 2, 0, 1), \qquad y, z \in \mathbb{R}.\]
    The solution $(0, 2, -8, 5)$ is the case $y = -8$, $z = 5$; it is not a multiple of $(-5, 1, 1, 0)$. The dimension is $2$.
    \item[\ref{pr:inverse2}.] $A^{-1} = \begin{pmatrix} 2 & -1 \\ -5 & 3 \end{pmatrix}$.
    \item[\ref{pr:inverse3}.] $A^{-1} = \begin{pmatrix} 0 & -1 & 1 \\ 1 & 0 & -1 \\ -1 & 1 & 1 \end{pmatrix}$ and $x = (1, -2, 4)^\top$.
    \item[\ref{pr:arithmetic}.] (a) $x = (1, -2, 1)^\top$. (b) row $1$ $-$ $2\cdot$row $2$ $+$ row $3$ $= 0$, so row operations produce a zero row.
    \item[\ref{pr:elementary}.] Swap the rows, multiply row $2$ by $\frac{1}{2}$, subtract $3$ times row $2$ from row $1$: $E_1 = \begin{pmatrix} 0 & 1 \\ 1 & 0 \end{pmatrix}$, $E_2 = \begin{pmatrix} 1 & 0 \\ 0 & \frac{1}{2} \end{pmatrix}$, $E_3 = \begin{pmatrix} 1 & -3 \\ 0 & 1 \end{pmatrix}$. Then $A^{-1} = E_3 E_2 E_1 = \begin{pmatrix} -\frac{3}{2} & 1 \\ \frac{1}{2} & 0 \end{pmatrix}$ and $A = E_1^{-1} E_2^{-1} E_3^{-1}$.
    \item[\ref{pr:idempotent}.] (a) $A = A^{-1}A^2 = A^{-1}A = I$. (b) $\ker(T_P) = \Span\{(1, -1)\}$ and $\im(T_P)$ is the $x$-axis: $T_P$ projects onto the $x$-axis along lines parallel to $(1, -1)$. (d) If $v = Aw$ and $Av = 0$, then $v = Aw = A^2 w = Av = 0$.
    \item[\ref{pr:nilpotent}.] $N^2 = \begin{pmatrix} 0 & 0 & 3 \\ 0 & 0 & 0 \\ 0 & 0 & 0 \end{pmatrix}$, $N^3 = 0$, and $(I - N)^{-1} = I + N + N^2 = \begin{pmatrix} 1 & 1 & 5 \\ 0 & 1 & 3 \\ 0 & 0 & 1 \end{pmatrix}$.
    \item[\ref{pr:interpolation}.] (a) $\begin{pmatrix} 1 & -1 & 1 \\ 1 & 0 & 0 \\ 1 & 1 & 1 \end{pmatrix}$, with determinant $2$. (c) $p = 1 - x^2$ and $q = 2 + 2x + x^2$.
    \item[\ref{pr:dims-special}.] (a) $3$, basis $e_{11}, e_{22}, e_{12} + e_{21}$; (b) $1$, basis $e_{12} - e_{21}$; (c) $6$; (d) $10$; (e) $6$; (f) $5$.
    \item[\ref{pr:groups}.] (a) Yes. (b) No: $\begin{pmatrix} 1 & 0 \\ 0 & 2 \end{pmatrix}\begin{pmatrix} 0 & 1 \\ 1 & 0 \end{pmatrix} = \begin{pmatrix} 0 & 1 \\ 2 & 0 \end{pmatrix}$ is not symmetric. (c) No: $\begin{pmatrix} 1 & 0 \\ 0 & -1 \end{pmatrix}$ and $\begin{pmatrix} 0 & 1 \\ 1 & 0 \end{pmatrix}$ square to $I_2$, but their product squares to $-I_2$. (d) Yes. (e) Yes; there are $n!$ of them, and each one is orthogonal.
    \item[\ref{pr:inverse-triangular}.] $\begin{pmatrix} 1 & -2 & 5 \\ 0 & 1 & -4 \\ 0 & 0 & 1 \end{pmatrix}$. For an upper triangular $A$ with nonzero diagonal entries, the Gauss--Jordan method only needs to scale rows and to subtract multiples of lower rows from higher rows, and these operations keep the right block upper triangular.
    \item[\ref{pr:orthogonal2}.] $Q = R_\theta$ (rotations, determinant $1$) or $Q = R_\theta \begin{pmatrix} 1 & 0 \\ 0 & -1 \end{pmatrix} = \begin{pmatrix} \cos\theta & \sin\theta \\ \sin\theta & -\cos\theta \end{pmatrix}$ (reflections, determinant $-1$).
    \item[\ref{pr:tridiagonal}.] (a) $D_1 = 2$, $D_2 = 3$, $D_3 = 4$. (c) $D_n = n + 1$.
    \item[\ref{pr:det4}.] $\det(A) = -1$, so $A$ is invertible.
    \item[\ref{pr:det-rules}.] $-6$, $18$, $54$, $-\frac{3}{2}$, $2$; and $\det(C) = -2$.
    \item[\ref{pr:vandermonde}.] (a) Subtract row $1$ from rows $2$ and $3$, and take out the factors $b - a$ and $c - a$. (b) The determinant is nonzero, so the map $p \mapsto (p(a), p(b), p(c))$ is an isomorphism.
    \item[\ref{pr:lambda}.] (a) $\lambda = 2$ and $\lambda = 5$; $\det(A - 3I_3) = -16 \neq 0$. (b) $x = (-8, 1, 1)$ for $\lambda = 2$, and $x = (-1, -1, 1)$ for $\lambda = 5$. (c) Every row sums to $11$, so $A(1, 1, 1)^\top = 11 (1, 1, 1)^\top$.
    \item[\ref{pr:upper5}.] (a) $\det(A) = 1 \cdot 2 \cdot 3 \cdot 4 \cdot 5 = 120$; yes; yes. (b) $(0, 0, 0, 0, 25)$. (c) $A - I_5$ and $A - 2I_5$ are triangular with a $0$ on the diagonal, so their determinants are $0$. Take $x = e_1$ and $y = (2, 1, 0, 0, 0)^\top$.
    \item[\ref{pr:area}.] (a) $\frac{1}{2}\left|\det \begin{pmatrix} 3 & 1 \\ 1 & 4 \end{pmatrix}\right| = \frac{11}{2}$. (b) $\det = 7$: the areas are $7$ and $7\pi$. (d) Only the reflection, with determinant $-1$. The rotation by $180^\circ$ has determinant $+1$.
    \item[\ref{pr:fibonacci}.] $M^{10} = \begin{pmatrix} 89 & 55 \\ 55 & 34 \end{pmatrix}$, and $89 \cdot 34 - 55^2 = 1$. Cassini follows from $\det(M^n) = \det(M)^n = (-1)^n$.
    \item[\ref{pr:trace}.] (a) For example, $A = \begin{pmatrix} 1 & 1 \\ 0 & 1 \end{pmatrix}$ and $B = \begin{pmatrix} 0 & 1 \\ 1 & 0 \end{pmatrix}$. $\tr(AB - BA) = 0 \neq 2 = \tr(I_2)$. (b) $\tr(A^\top A)$ is the sum of the squares of all entries. (c) $\tr(Ae_{ji}) = A_{ij}$.
    \item[\ref{pr:01-matrices}.] (a) $6$. (For $3 \times 3$ matrices with entries in $\{0, 1\}$, a computer finds that $174$ of the $512$ are invertible.) (b) If $A$ and $A^{-1}$ have integer entries, then $\det(A)\det(A^{-1}) = 1$ is a product of two integers. Conversely, use the cofactor formula $A^{-1} = \frac{1}{\det(A)} C^\top$.
    \item[\ref{pr:bases-B}.] $\RREF(B)$ has rows $(1, 2, 0, 3)$, $(0, 0, 1, 2)$, $(0, 0, 0, 0)$, and $\rank(B) = 2$. Bases: $\Null(B)$: $(-2, 1, 0, 0)$, $(-3, 0, -2, 1)$; $\col(B)$: $(1, 2, 1)$, $(0, 1, 1)$; $\row(B)$: $(1, 2, 0, 3)$, $(0, 0, 1, 2)$. Indeed, $2 + 2 = 4$.
    \item[\ref{pr:onto}.] (a) $5 - 3 = 2$. (b) Not surjective, since the rank is at most $3$; it can be injective. (c) No: $\dim \Null(A) \geq 6 - 4 = 2$. (d) $3$.
    \item[\ref{pr:intersect}.] (a) The $7 \times 8$ matrix has a nonzero null vector $(a_1, \ldots, a_4, b_1, \ldots, b_4)$; then $\sum a_i v_i = -\sum b_j w_j$ is a nonzero vector in $V \cap W$. (c) $\Span\{e_1, e_2, e_3\}$ and $\Span\{e_4, e_5, e_6\}$.
    \item[\ref{pr:rankAB}.] (b) $\rank(AB) \leq \rank(A) \leq 2 < 3$.
    \item[\ref{pr:ATA}.] (a) $x^\top A^\top A x = \sum_i (Ax)_i^2$. (b) Use the rank--nullity theorem. (c) $\Null(A^\top A) = \Null(A) = \{0\}$.
    \item[\ref{pr:eig2}.] (a) $p_A(\lambda) = \lambda^2 - 3\lambda - 4 = (\lambda - 4)(\lambda + 1)$; $E_4 = \Span\{(2, 3)^\top\}$ and $E_{-1} = \Span\{(1, -1)^\top\}$. (b) $4 + (-1) = 3$ and $4 \cdot (-1) = -4$. (c) $P = \begin{pmatrix} 2 & 1 \\ 3 & -1 \end{pmatrix}$, $D = \diag(4, -1)$, $P^{-1} = \frac{1}{5}\begin{pmatrix} 1 & 1 \\ 3 & -2 \end{pmatrix}$. (d) $A^k = \frac{1}{5}\begin{pmatrix} 2 \cdot 4^k + 3(-1)^k & 2 \cdot 4^k - 2(-1)^k \\ 3 \cdot 4^k - 3(-1)^k & 3 \cdot 4^k + 2(-1)^k \end{pmatrix}$.
    \item[\ref{pr:eig-triangular}.] The eigenvalues are the diagonal entries $2$, $3$ and $1$. The eigenspaces are spanned by $(1, -1, 5)^\top$ (for $2$), $(0, -2, 1)^\top$ (for $3$) and $(0, 0, 1)^\top$ (for $1$). Diagonalizable, since the three eigenvalues are distinct.
    \item[\ref{pr:eig-defective}.] The eigenvalue $2$ has algebraic multiplicity $2$, but geometric multiplicity $3 - \rank(N - 2I_3) = 1$; the eigenvalue $3$ has both multiplicities equal to $1$. Not diagonalizable.
    \item[\ref{pr:eig-geometric}.] (a) $1$ on the $x$-axis, $-1$ on the $y$-axis. (b) $1$ on $\Span\{(1, 1)\}$, $0$ on $\Span\{(1, -1)\}$. (c) Only $1$, with eigenspace the $x$-axis. (d) No real eigenvalues; the complex eigenvalues are $\pm i$. (e) $2$ on the $x$-axis, $3$ on the $y$-axis. (f) Only $-1$, with eigenspace $\mathbb{R}^2$. All are diagonalizable except (c) and (d) (over $\mathbb{R}$).
    \item[\ref{pr:eig-gre}.] (a) Yes (three distinct eigenvalues); yes ($0$ is not an eigenvalue). (b) $5$, $-8$, $1 + 4 + 16 = 21$, $(-8)^3 = -512$. (c) $A^{-1}$: $-1, \frac{1}{2}, \frac{1}{4}$; $A^2$: $1, 4, 16$; $A + 2I_3$: $1, 4, 6$.
    \item[\ref{pr:eig-idempotent}.] (a) If $Pv = \lambda v$, then $\lambda v = P^2 v = \lambda^2 v$, so $\lambda^2 = \lambda$. (b) If $v = Pw$, then $Pv = P^2w = Pw = v$. Bases of $\im(T_P)$ and $\ker(T_P)$ together span $\mathbb{R}^n$ and are linearly independent, so they form a basis of eigenvectors. (c) Only ``$P$ is diagonalizable''; $P = \begin{pmatrix} 1 & 1 \\ 0 & 0 \end{pmatrix}$ is a counterexample to the other two.
    \item[\ref{pr:eig-props}.] (a) For the shear, $E_1 = \Span\{e_1\}$, but for its transpose, $E_1 = \Span\{e_2\}$. (c) $v^\top v = (Qv)^\top(Qv) = \lambda^2 v^\top v$. (d) $A^k v = \lambda^k v = 0$. (e) $P^{-1} I_2 P = I_2$ for every $P$.
    \item[\ref{pr:eig-recurrence}.] (a) $M = \begin{pmatrix} 1 & 2 \\ 1 & 0 \end{pmatrix}$. (b) Eigenvalues $2$ and $-1$, with eigenvectors $(2, 1)^\top$ and $(-1, 1)^\top$; $a_n = \frac{2^n - (-1)^n}{3}$. (c) $a_{10} = \frac{1024 - 1}{3} = 341$.
    \item[\ref{pr:eig-markov}.] (a) $\lambda = 1$ with eigenvector $(4, 3)^\top$, and $\lambda = 0.3$ with eigenvector $(1, -1)^\top$; the stationary distribution is $\big(\frac{4}{7}, \frac{3}{7}\big)^\top$. (b) $x_k = \big(\frac{4}{7}, \frac{3}{7}\big)^\top + \frac{3}{7}(0.3)^k (1, -1)^\top$; $x_2 = (0.61, 0.39)^\top$; $x_k \to \big(\frac{4}{7}, \frac{3}{7}\big)^\top$.
    \item[\ref{pr:eig-power}.] (a) $(2, 1)$, $(5, 4)$, $(14, 13)$, $(41, 40)$; the ratio tends to $1$, i.e., to the direction of the eigenvector $(1, 1)$. (b) $x_0 = \frac{1}{2}(1, 1)^\top + \frac{1}{2}(1, -1)^\top$, so $A^k x_0 = \frac{1}{2}(3^k + 1, \ 3^k - 1)^\top$.
    \item[\ref{pr:eig-complex}.] (a) $p(\lambda) = (a - \lambda)^2 + b^2$. (b) $3 \pm 4i$; the matrix is $5 R_\theta$ with $\cos\theta = \frac{3}{5}$ and $\sin\theta = \frac{4}{5}$: $r = |3 \pm 4i| = 5$, and $\theta$ is the angle of $3 + 4i$.
\end{description}

\end{document}
